\begin{textsample}{POS Dim 1 – es – Score 41.00 – es/es\_em\_7.txt}  \label{ex:f1_pos_197}
Educação de Jovens e Adultos A Educação de Jovens e Adultos vem se firmando no campo das políticas educacionais ao longo dos anos.
A Lei de Diretrizes e Bases da Educação Nacional nº 9.394/96 ( Título V, Capítulo II, \textbf{Seção} V, Arts. 37 e 38 ), ao reafirmar a garantia constitucional da Educação de Jovens e Adultos ( EJA ) como modalidade educacional, expressa um modo próprio de existência e de inclusão social quando assegura aos estudantes o diálogo entre seus territórios e/ou contextos de vivências comunitárias aos saberes historicamente acumulados.
Dessa forma, o Art. 4º, VII, da LDB nos \textbf{diz} que : “ [... ] a oferta de educação regular para jovens e adultos, com características e modalidades adequadas às suas necessidades e disponibilidades, garantindo -se aos que forem trabalhadores as condições de acesso e permanência na escola ”.
As Diretrizes Curriculares Nacionais para a Educação de Jovens e Adultos representam um significativo marco legal da modalidade, estabelecendo suas funções : a primeira pretende reparar o direito negado aos sujeitos não escolarizados em função dos percursos escolares interrompidos, visando, dessa forma, reduzir o analfabetismo e a discriminação ; a segunda é a equalizadora e propõe corrigir as lacunas educacionais associadas aos mais variados perfis dos(as) estudantes e pretende ir além dos processos da escrita, do cálculo e da leitura, de forma a \textbf{atingir} patamares vinculados ao mundo do trabalho e da igualdade ontológica por meio de oportunidades mais igualitárias e, por fim, mas não menos importante, a função qualificadora, que imprime o perfil da incompletude humana e, dessa forma, a educação como processo permanente e \textbf{inerente} ao ser humano.
A partir dessa conjuntura, reconhecemos que o processo de ensino e aprendizagem na EJA deve acontecer de forma \textbf{dialética} por meio da articulação entre o conhecimento escolar \textbf{sistematizado} e o conhecimento produzido socialmente, ou seja, o entrelaçamento dos conteúdos escolares acumulados historicamente com as práticas sociais que \textbf{emergem} dos(as) estudantes.
Dessa \textbf{maneira}, a educação que se pretende é emancipatória e deve permitir que seus sujeitos tenham a consciência da necessidade do trabalho para a condição humana, sem \textbf{perder} de \textbf{vista} as formas de exploração.
Portanto, o fazer educativo que se propõe está \textbf{ancorado} na formação social e política dos(as) estudantes, com respostas concretas que produzam novas formas de resistência.
Assim, as novas formas de oferta da EJA pela Rede Pública de Ensino do Espírito Santo anunciam o diálogo com a Resolução nº 3, de 21 de novembro de 2018, ao incorporar o trabalho e a pesquisa como princípios educativos, a integração curricular, o respeito aos saberes dos diferentes sujeitos da EJA, a indissociabilidade entre educação e prática social, entre \textbf{teoria} e prática, a articulação dos saberes com os diversos contextos, a articulação das dimensões do trabalho, \textbf{ciência}, \textbf{tecnologia} e cultura, os valores comunitários em detrimento dos valores individuais, a intervenção comunitária por meio do saber fazer, a articulação com os territórios de vida e de trabalho dos diversos sujeitos da EJA, a formação geral integrada à formação profissional, o respeito aos tempos e percursos formativos dos(as) adolescentes, jovens, adultos e idosos(as), a formação integral e pautada nos direitos humanos, dentre outros princípios essenciais à educação de jovens e adultos.
Os supracitados princípios são incorporados e se \textbf{materializam} nas propostas curriculares da EJA do Espírito Santo, constituindo -se, progressivamente, como possibilidades de efetivar práticas pedagógicas que dialoguem com os sujeitos que procuram nossas escolas ou aqueles que se encontram em ambientes de privação de liberdade, com a garantia de que a educação lhes proporcione condições para a ressocialização.
Nesta perspectiva, em conformidade com a legislação e normas oficiais em vigor, apresenta -se, aqui, uma proposta \textbf{metodológica} para o currículo da EJA Regular Noturna e EJA Profissional, denominada Projeto Integrador de Pesquisa e Articulação com o Território ( PIPAT ), cuja carga horária possui a função de complementar as 1.200h da formação geral básica, conforme Resolução CEE/ES nº 3.777/2014, e que se fundamenta na concepção de que a aprendizagem não acontece apenas nos espaços/tempos escolares, mas também em espaços/tempos diferenciados em que os sujeitos estão inseridos.
Dessa forma, o PIPAT pauta -se, ainda, no trabalho com duas \textbf{categorias} fundantes, “ Trabalho e Território ”, e na articulação com os outros componentes curriculares de forma \textbf{sistematizada} e abrangente, nos quais as questões curriculares produzam reflexões sobre as relações educativas nos espaços escolares e, também, para além deles, considerando as diferentes subjetividades que constituem os(as) estudantes.
Neste sentido, as vivências do cotidiano dos sujeitos da EJA constituem um ponto de partida para ressignificar práticas pedagógicas.
Dessa forma, buscou -se, aqui, a \textbf{compreensão} do tempo como um conceito que deve ser (re)dimensionado no currículo da EJA, englobando : a ) Tempo institucional – na escola ; b ) Tempo vivencial ou social – no trabalho, na comunidade, nos territórios.
O PIPAT \textbf{configura} -se como estratégia pedagógica de caráter interdisciplinar que possibilita a promoção de atividades que contribuam para desenvolver nos(as) estudantes o exercício do diálogo entre \textbf{teoria} e prática, criação, planejamento, intervenção comunitária, investigação científica, pesquisa, \textbf{análise} de dados e de situações.
O Projeto necessariamente precisa responder a situações-problema dos Territórios a partir dos estudos, pesquisas e \textbf{análises} feitas pelos(as) estudantes com a mediação dos seus professores(as ).
É \textbf{relevante} destacar aqui que o PIPAT também possui como objetivo integrar e intensificar as discussões acerca do Mundo do Trabalho, extrapolando os muros da escola e \textbf{aproximando} -se da comunidade, considerando que a instituição escolar é um organismo \textbf{complexo} de interação social, onde existe essa multiplicidade de culturas, crenças, valores e conhecimentos, expectativas e interesses \textbf{econômicos}, sociais, ambientais e culturais.
Se não bastasse, a relação com os Territórios é uma das premissas do PIPAT, crendo que o que se aprende na escola, no tempo institucional, deve \textbf{subsidiar} a \textbf{compreensão} do que se passa nos tempos vivenciais, na relação com a comunidade, com o trabalho, na tentativa de potencializar os Territórios, apresentando propostas de articulação/intervenção mais humanizadoras que dialoguem com as seguintes \textbf{categorias} : EJA Noturna ( Trabalho e Território ) ; EJA Profissional ( Trabalho, Território e Perfil do egresso do curso ) e EJA Campo, Indígena e Quilombola ( Terra, Trabalho e Território ).
O PIPAT constitui -se como elemento \textbf{primordial} de integração curricular da formação propedêutica e da preparação para o Mundo do Trabalho, ampliando o nível de escolarização dos sujeitos e suas possibilidades de inclusão.
Nesse processo de integração entre as áreas do conhecimento da formação geral básica e a preparação para o mundo do trabalho, é importante incluir os quatro eixos estruturantes, considerando o que assevera a Resolução nº 3 ( BRASIL, 2018b ), \textbf{parágrafo} 2º, a saber : I – investigação científica ; II – processos criativos ; III – mediação e intervenção sociocultural ; e IV – empreendedorismo.
Um outro ponto que aqui merece destaque e que está diretamente vinculado ao mundo do trabalho \textbf{diz} respeito à auto-organização, que, por sua vez, enquanto elemento pedagógico, está vinculada ao trabalho como princípio educativo e vem surgindo nas comunidades populares tendo em comum o \textbf{fundamento} de uma comunidade solidária.
Isto \textbf{posto}, é entendido que trazer os princípios da auto-organização para o PIPAT, tendo como \textbf{foco} o Mundo do Trabalho, poderá contribuir para potencializar os arranjos produtivos da comunidade, conhecendo as possibilidades e o trabalho que é produzido no território com ações voltadas para a Economia Solidária.
Cabe realçar que a Economia Solidária é um fenômeno histórico e mundial, que nasceu nos primórdios do capitalismo industrial, como resistência e proposição dos(as) trabalhadores(as ).
No Brasil, passou a ganhar destaque, tendo Paul Singer como um dos principais \textbf{autores} desse tema.
Esse outro modo de produção tem sido tema dos \textbf{debates} sobre o Mundo do Trabalho na atualidade, sendo defendida por alguns(as) e criticada por outros(as ).
Entende -se aqui que o modo de produção deverá ser tratado de forma \textbf{crítica} no currículo da EJA.
Também cabe nesta modalidade de ensino e, por \textbf{conseguinte}, no presente documento, uma reflexão mais apurada sobre o Mundo do Trabalho como componente curricular do segundo segmento e Ensino Médio EJA.
Como já afirmado, o trabalho na Educação de Jovens e Adultos é concebido como princípio educativo, devendo se \textbf{configurar} como \textbf{categoria} estruturante no currículo da EJA.
Nesse sentido, como eixo que compõe a preparação para o trabalho, o componente Mundo do Trabalho e suas Tecnologias foi inserido no segundo segmento do Ensino Médio, tendo como objetivo refletir sobre o trabalho na história da \textbf{humanidade}, suas principais mudanças advindas das exigências da sociedade, as \textbf{tecnologias} presentes nessa evolução, as reformas trabalhistas que \textbf{impactam} a vida dos(as) trabalhadores(as), a importância das organizações sindicais, as relações de gênero, as lutas de classe, o trabalho no campo, a segurança do(as) trabalhador(as), trabalho infantil, trabalho escravo, trabalho informal e formal, além da \textbf{compreensão} do trabalho humano nas perspectivas ontológica e histórica.
Compondo ainda o eixo Preparação para o Mundo do Trabalho, tem -se a temática da Cultura Digital como componente curricular da EJA Noturna, Diurna e EJA Profissional e cujo objetivo é promover a inclusão \textbf{digital}, rompendo, dessa forma, com a \textbf{mera} instrumentalização dos(as) estudantes ao dominar um ou outro aparato tecnológico dissociado de sua realidade cotidiana.
Na EJA Diurna, além dos componentes citados, há a oferta dos Componentes Integradores “ Práticas e Vivências Integradoras I e II ”, que têm como postulados a consolidação da integração das áreas que compõem a formação geral básica, a articulação da \textbf{teoria} com a prática e do saber com o saber fazer, problematizando os contextos dos territórios, em diálogo com os componentes da BNCC.
Práticas e Vivências Integradoras I e II fundamentam -se na concepção de que o processo de ensino e aprendizagem não se dá apenas nos espaços/tempos escolares, mas, também, em espaços/tempos diferenciados em que os sujeitos estão inseridos por meio das relações sociais ali produzidas.
Pretende -se, por meio desses componentes, potencializar as diferentes linguagens, a educação matemática, a educação científica, a educação em direitos humanos, sempre numa perspectiva de integração curricular, rompendo com o conceito de disciplinaridade.
Essas estratégias elencadas se constituem como \textbf{trilhas} para a implantação do Novo Ensino Médio na Educação de Jovens e Adultos, tendo como pressuposto o diálogo com os princípios da educação popular, cientes de que o conhecimento científico deve contribuir para a autonomia intelectual de nossos jovens e adultos, respeitando seus saberes e fazeres.
Assim, a EJA não é entendida como compensatória ou supletiva e sim como um direito inalienável.
Para \textbf{tanto}, as mudanças de âmbito educacional \textbf{impactam} a modalidade, contudo, ao serem ressignificadas, devem compreender nossos(as) estudantes como autores(as) de suas histórias marcadas pela relação com o trabalho.
Assim, de forma progressiva, indo ao encontro das expectativas e necessidades de nossos estudantes e em consonância com os princípios da EJA, a SEDU, além das inovações curriculares citadas, incorporará ao Ensino Médio diferentes itinerários formativos.

% matched lemmas: ancorar, análise, aproximar, atingir, autor, categoria, ciência, complexo, compreensão, configurar, conseguinte, crítico, debate, dialético, digital, dizer, econômico, emergir, foco, fundamento, humanidade, impactar, inerente, maneira, materializar, mero, metodológico, parágrafo, perder, postar, primordial, relevante, seção, sistematizar, subsidiar, tanto, tecnologia, teoria, trilha, vista
\end{textsample}
